(1) 令 \(b=\beta\hbar\omega\)、\(r=e^{-b}\)。能级 \(E_n=\hbar\omega(n+\tfrac12)\)，配分函数 \(Z=e^{-b/2}/(1-r)\)。因此
\[
\rho=(1-r)\sum_{n=0}^\infty r^n|n\rangle\langle n|,\quad
\bar n=\frac r{1-r}=\frac1{e^b-1},\quad \langle H\rangle=\hbar\omega(\bar n+\tfrac12).
\]
(2) 用 \(x=\sqrt{\hbar/(2m\omega)}(a+a^\dagger)\)、\(p=-i\sqrt{m\hbar\omega/2}(a-a^\dagger)\)，热态对角性使奇矩为零：
\[
\langle x^2\rangle=\frac{\hbar}{2m\omega}\coth\frac b2,\quad
\langle p^2\rangle=\frac{m\hbar\omega}{2}\coth\frac b2.
\]
故 \(\Delta x\Delta p=\frac\hbar2\coth(b/2)\ge\hbar/2\)，零温取等，高温给经典均分 \(\langle x^2\rangle\simeq k_BT/(m\omega^2)\)、\(\langle p^2\rangle\simeq mk_BT\)。
(3) 热态是零均值Gaussian态，令 \(t=\tanh(b/2)\)，其Wigner函数为
\[
W(x,p)=\frac{t}{\pi\hbar}\exp\left[-\frac{2t}{\hbar\omega}\left(\frac{p^2}{2m}+\frac{m\omega^2x^2}{2}\right)\right].
\]
两个高斯积分分别给 \(\int W\,dx\,dp=1\) 与 \(\int x^2W\,dx\,dp=\hbar/(2m\omega t)\)、\(\int p^2W\,dx\,dp=m\hbar\omega/(2t)\)，与第二问相同。其非负性源于热Gaussian态，不代表任意量子态Wigner函数非负。
(4) 几何本征值 \(p_n=(1-r)r^n\) 给
\[
\operatorname{Tr}\rho^2=(1-r)^2/(1-r^2)=\frac{1-r}{1+r}=\tanh(b/2),
\]
\[
S=k_B[(\bar n+1)\ln(\bar n+1)-\bar n\ln\bar n].
\]
\(T\to0\) 时 \(\bar n\to0\)、纯度趋1、熵趋0；\(k_BT\gg\hbar\omega\) 时 \(\bar n\simeq k_BT/(\hbar\omega)\)、纯度 \(\simeq\hbar\omega/(2k_BT)\)、熵 \(S\simeq k_B[1+\ln(k_BT/\hbar\omega)]\)，与经典相空间热熵一致。
